

%-------------------- Fig --------------------
\begin{figure}[t]
\centering
\includegraphics[width=1.0\columnwidth]{fig_dpl.eps}
\caption{
Overview of the free energy principle~\cite{Friston2010nrn}.
}
\label{fig_dpl}
\hspace{10mm}
\end{figure}



\section{Deep Predictive Learning}

%This section explains the four strategies that form the basis of the deep predictive learning (DPL) concept, along with the free energy principle and predictive coding. Then, as an implementation of the DPL concept, we describe the phases for building a basic predictive model and three particularly important elements in that phase. Also, we discuss the concept from the perspective of embodiment~\cite{Brooks1991intelligence} and dynamic cognitive system~\cite{Tani1996model}.
This section explains the four strategies that form the basis of the deep predictive learning (DPL) concept, free-energy principle, and predictive coding. Then, as an implementation of the DPL concept, we describe the phases for building a basic predictive model and the three particularly important elements in that phase. In addition, we discussed this concept from the perspectives of embodiment~\cite{Brooks1991intelligence} and dynamic cognitive systems~\cite{Tani1996model}.



\subsection{Main Concept and Policies}

%The Deep Neural Network (DNN) framework uses large amounts of data to optimize models, and the goal is for that model to function appropriately even with new data. However, it is usually not easy to obtain an optimal model because it is impossible to know which state will provide the greatest generalizability. Especially in robots, it can be assumed that they are constantly exposed to unexpected situations because they operate in the real environment. Therefore, an approach is required that assumes that the learning is incomplete and considers how to deal with prediction errors.
The Deep Neural Network (DNN) framework uses large amounts of data to optimize models, and the goal is for the model to function appropriately, even with new data. However, obtaining an optimal model is usually difficult because it is impossible to determine the state that provides the greatest generalizability. It can be assumed that robots are constantly exposed to unexpected situations because they operate in real environments. Therefore, an approach that assumes that learning is incomplete and considers how to manage prediction errors is required.


%Here, there is an information theory of the brain called the free energy principle~\cite{Friston2006jpp,Friston2010nrn} proposed by K. Friston that can help tackle the above problem. The free energy principle states that the perception, learning, and behavior of living things are determined by minimizing a cost function called variational free energy (Figure~\ref{fig_dpl}). We try to summarize it below.
There is an information theory of the brain called the free energy principle~\cite{Friston2006jpp,Friston2010nrn} proposed by K. Friston, which can help address the above problem. The free-energy principle states that living things' perception, learning, and behavior are determined by minimizing a cost function called variational free energy (Figure~\ref{fig_dpl}). We summarize this as follows:
\begin{equation}
    FreeEnergy = -\log p(x|z) + D_{KL}[q(z|x)||p(z)] \label{eq_free}
\end{equation}
%The above corresponds to the negative Evidence Lower Bound in general machine learning models, and this equation itself is not new. However, by considering the interpretation of the equation and the multiple parameters involved, we can discuss the design of a system with a concrete body. The brain is a predictor and has a model $p(z)$ as a priori belief about the perception $z$ from the environment that causes the sensory input $x$. In addition, perception $z$ includes not only the hidden states of exteroceptive sensations such as vision and touch but also the hidden states of proprioception (e.g., joints) and interoceptive sensations related to kinesthetic and visceral sensations. The brain generates a sensory prediction $g_w(z)$ based on this perception $z$, and obtains a posteriori belief $q(z|x)$ by receiving the current interoceptive sensory input $x$. After that, the brain continuously generates predictions of the next sensory input based on the updated posterior beliefs and updates the beliefs. The predicted sensory inputs include a prediction of accuracy (the reciprocal of the sensation variance $v_x$). In other words, the first term in the equation \ref{eq_free} refers to the prediction error weighted by this accuracy, and the second term refers to the influence of the complexity of the structure of the prior belief (model).
The above equation corresponds to the negative Evidence Lower Bound in general machine learning models, which is not new. However, by considering the interpretation of the equation and the multiple parameters involved, we can discuss the design of a system with a concrete body. The brain is a predictor with a model $p(z)$ as an a priori belief about the perception $z$ from the environment that causes sensory input $x$. In addition, perception $z$ includes not only the hidden states of exteroceptive sensations, such as vision and touch but also the hidden states of proprioception (e.g., joints) and interoceptive sensations related to kinesthetic and visceral sensations. The brain generates a sensory prediction $g_w(z)$ based on this perception $z$ and obtains a posteriori belief $q(z|x)$ by receiving the current interoceptive sensory input $x$. Subsequently, the brain continuously generates predictions for the next sensory input based on updated posterior beliefs and updates those beliefs. The predicted sensory inputs included the prediction accuracy ( reciprocal of the sensation variance $v_x$). In other words, the first term in equation \ref{eq_free} refers to the prediction error weighted by this accuracy. The second term refers to the influence of the complexity of the prior belief (model) structure.


%Strategies for reducing prediction errors are called "predictive coding," and can be divided into four main types. By combining these strategies, the DPL realizes online motion generation tailored to the environment.
Strategies for reducing prediction errors are called "predictive coding" and can be classified into four main types. The DPL realizes online motion generation that is tailored to the environment by combining these strategies.


\begin{description}
\item[Training]\mbox{}\\
%The basic strategy is to optimize the entire model. The model that generates sensory predictions is trained, and the internal weight parameter $w$ is changed according to the prediction error. However, training the DNN model is costly and difficult to implement under time constraints. This problem is particularly serious in deep learning approaches, as the model size tends to be large. Furthermore, depending on the real robot situation, frequent model modifications may be inappropriate.
The basic strategy is to optimize the entire model. The model that generates the sensory predictions is trained, and the internal weight parameter $w$ changes based on the prediction error. However, training a DNN model is expensive and difficult under time constraints. This problem is particularly serious in deep learning approaches because the model size tends to be large. Furthermore, frequent model modifications may be inappropriate depending on the robot's situation.


\item[Modification of Posterior Belief]\mbox{}\\
%The second strategy is to modify the hypothetical posterior belief $q(z)$. The values of sensory input are modified in a closed manner based on the predicted perception, and as a result, the prediction error of the model is reduced. For example, when humans perceive the environment, they often interpret the real world as an afterthought based on previous images in their brains. And, the brain is unaware of the errors. This strategy is a method of adapting to changes in the real world by using predictive biases to create a virtual perception that differs from reality (Section~\ref{sec_generation}).
The second strategy is to modify hypothetical posterior belief $q(z)$. The values of the sensory input were modified in a closed manner based on the predicted perception; as a result, the prediction error of the model was reduced. For example, when humans perceive an environment, they often interpret the real world as an afterthought based on previous brain images. Additionally, the brain is unaware of these errors. This strategy is a method of adapting to real-world changes using predictive biases to create a virtual perception that differs from reality (Section~\ref{sec_generation}).


\item[Active Inference]\mbox{}\\
%The third strategy is to use the body to influence the external environment and change the sensory input $x$ to reduce prediction errors. This operation is called active inference in the context of the free energy principle. There are methods such as manipulating the environment by oneself or changing the viewpoint by changing one's position. This strategy is essential for autonomous agents that have bodies in the real world, such as robots. We implement the behavioral transitions for performing posterior beliefs and active inference using attractor dynamics embedded in the internal state of the model (Sectoin~\ref{sec_imple}).
The third strategy is to use the body to influence the external environment and change the sensory input $x$ to reduce the prediction errors. This operation is called active inference in the context of the free energy principle. These methods include manipulating the environment by oneself or changing viewpoints by changing one's position. This strategy is essential for autonomous agents, such as robots, with bodies in the real world. We implemented behavioral transitions to perform posterior belief and active inference using attractor dynamics embedded in the model's internal state (Sectoin~\ref{sec_imple}).


\item[Modification of Prediction Accuracy]\mbox{}\\
%The fourth strategy is to reduce the apparent error by setting the prediction accuracy $1/v_x$ of sensory input $x$ small and multiplying it by the weight of the prediction error. This operation, which increases the precision of sensory input, is associated with directing attention. From the perspective of predictive encoding, this strategy can be interpreted as placing emphasis on prediction errors and increasing prediction corrections when sensory input is reliable (when attention is directed). On the other hand, it can be interpreted that when sensory input is unreliable (when attention is not directed), one ignores the prediction error and maintains one's inference. The implementation example corresponds to Section~\ref{sec_attention}.
The fourth strategy involves reducing the apparent error by setting the prediction accuracy $1/v_x$ of the sensory input $x$ and multiplying it by the weight of the prediction error. This operation, which increases the precision of sensory input, is associated with directing attention. From the perspective of predictive encoding, this strategy can be interpreted as emphasizing prediction errors and increasing prediction corrections when sensory input is reliable (when attention is directed). However, when sensory input is unreliable (when attention is not directed), the prediction error is ignored, and the inference is maintained. The implementation example corresponds to that presented in Section~\ref{sec_attention}.
\end{description}


%-------------------- Fig --------------------
\begin{figure*}[t]
\centering
\includegraphics[width=17cm]{1_overview.eps}
\caption{
%An example of a deep predictive learning framework, consisting of a feature extraction unit and a time-series prediction unit.
An example of a deep predictive learning framework consisting of a feature extraction unit and a time-series prediction unit.
}
\label{fig_overview}
\end{figure*}



\subsection{Implementation}

\subsubsection{Learning Sensorimotor Dynamics}
\label{sec_imple}

%We explain the construction phase of the DPL model to realize the above-mentioned strategies, divided into training data collection and model learning. Figure~\ref{fig_overview} shows the overall flow including the generation phase described in section \ref{sec_generation}.
We explain the construction phase of the DPL model to realize the strategies mentioned above, which are divided into training data collection and model learning. Figure~\ref{fig_overview} illustrates the overall flow, including the generation phase described in Section \ref{sec_generation}.


\textbf{Collecting Training Dataset:}
%In the DPL concept, a model to generate robot motions is acquired in a data-driven manner by directly using sensorimotor information from the robot (time-series data consisting of motion information, sensor information, etc.) as training data. Therefore, the training data must include information about the interaction between the robot's body and the environment. Furthermore, since the quantity and quality of training data greatly affect model performance, it is desirable to be able to efficiently collect high-quality demonstration data.
In the DPL concept, a model for generating robot motions is acquired in a data-driven manner by directly using sensorimotor information from the robot (time-series data consisting of motion and sensor information, etc.) as training data. Therefore, the training data must include information on the robot's body and environment interaction. Furthermore, because the quantity and quality of training data greatly affect model performance, it is desirable to collect high-quality demonstration data efficiently.


%In the first step, training data is collected by actually having the robot perform the desired task and recording motion information such as joint angles, and sensor information such as camera images at a fixed sampling rate. Typical teaching methods for robot motion include programming description, direct teaching, and teleoperation. Describing the motion in advance using a robot programming language is the simplest teaching method, but the description becomes complicated for long-horizon tasks. On the other hand, by having a human remotely control the robot and teaching it its motions, it is possible to obtain training data without precise modeling or parameter adjustment. Especially, the method of remotely controlling the manipulator from the actual robot's perspective (Wizard of Oz method) is desirable because it allows the robot to intuitively teach human operation skills for tasks. By controlling the robot as if they were controlling their own bodies, workers can interact naturally with the environment. The acquired teaching data are expected to include the information necessary for motion learning.
In the first step, training data are collected by having the robot perform the desired task and record motion information, such as joint angles, and sensor information, such as camera images, at a fixed sampling rate. Typical teaching methods for robot motion include programming, direct teaching, and teleoperation. Describing the motion in advance using a robot programming language is the simplest teaching method; however, the description becomes complicated for long-horizon tasks. However, training data can be obtained without precise modeling or parameter adjustment by having a human remotely control the robot and teaching its motion. In particular, remotely controlling the manipulator from an actual robot's perspective (the Wizard of Oz method) is desirable because it allows the robot to intuitively teach human operational skills for tasks. Workers could interact naturally with the environment by controlling the robot as if it were controlling its own body. The acquired teaching data were expected to include the information necessary for motion learning.


\textbf{Integration Learning of Sensorimotor Dynamics:}
%The learning target of our DPL concept is the time-series relationship between sensorimotor information in a system where the environment and the body dynamically interact. The training data does not have ground truth labels, and the model is trained to predict the next step robot state by taking the current robot state as input. This autoregressive learning eliminates the need for detailed design of a physical model of the environment, which is required in conventional robotics. In addition, the model can be expressed as dynamics that integrate environmental recognition and motion generation functions that span multiple modalities.
The learning target of our DPL concept is the time-series relationship between sensorimotor information in a system in which the environment and body dynamically interact. The training data did not contain ground truth labels, and the model was trained to predict the next-step robot state using the current robot state as the input. This autoregressive learning eliminates the need for a detailed design of the physical model of the environment, which is required in conventional robotics. In addition, the model can be expressed in terms of dynamics that integrate the environmental recognition and motion generation functions spanning multiple modalities.


%We construct a sensory and motor prediction system by combining multiple DNN models corresponding to the feature extraction part and the time-series prediction unit. The feature extraction unit extracts features from the sensor values acquired by the robot, and the time-series prediction unit learns sensorimotor information that integrates the extracted features and the robot state expressed by joint angles. We published the first model in 2014~\cite{Noda2014multimodal}, and since then this study has continued to expand. Each part can be connected end-to-end or learned independently. In this Section, we will clearly explain the role of each part.
We constructed a sensory and motor prediction system by combining multiple DNN models corresponding to the feature extraction and time-series prediction units. The feature extraction unit extracts features from the sensor values acquired by the robot, and the time-series prediction unit learns sensorimotor information that integrates the extracted features and the robot state expressed by joint angles. We published the first model in 2014~\cite{Noda2014multimodal}, and since then, this study has continued to expand. Each part can be connected end to end or learned independently. In this section, we explain the role of each component.


\begin{description}
\item[Feature Extraction Unit:]\mbox{}\\
%Sensor inputs that robots obtain from the real environment, such as cameras and tactile sensors, are often high-dimensional, so appropriate feature extraction is important. To appropriately learn sensorimotor information, the extracted sensor features must reflect the relationship between the robot body and the manipulated object. This is because it affects the prediction performance of motion information in the time-series learning section, such as in the case of approaching a manipulated object. Therefore, to obtain as detailed information as possible, it is desirable that the feature extractor can handle high-dimensional sensor input as is.
Sensor inputs that robots obtain from a real environment, such as cameras and tactile sensors, are often high dimensional; therefore, appropriate feature extraction is important. Sensorimotor information is appropriately learned by extracting sensor features that reflect the relationship between the robot body and the manipulated object. The extraction of sensor features affects the prediction performance of the motion information in the time-series learning section, such as when approaching a manipulated object. Therefore, to obtain as detailed information as possible, the feature extractor should handle high-dimensional sensor inputs as is.


%Here, as an easy-to-understand example, we consider image data. Raw image data contains a lot of information that is unnecessary for robot tasks, such as background, and it can become noise for training models. Many of the models we used adopted a convolutional auto-encoder (CAE) as a feature extractor. In the CAE, low-dimensional features are acquired in the middle layer by performing self-supervised learning to provide an output value that is equal to the input value. The encoded feature vector provides high-dimensional input information using fewer dimensions, making it robust against environmental noise. Furthermore, the above feature vectors can be narrowed down to information that is even more effective for task learning by introducing the attention mechanism (Section~\ref{sec_attention}).
Here, as an easy-to-understand example, we consider image data. Raw image data contain much unnecessary information for robot tasks, such as the background, and can become noise for training models. Many models adopted a convolutional autoencoder (CAE) as the feature extractor. In the CAE, low-dimensional features are acquired in the middle layer by performing self-supervised learning to provide an output value equal to the input value. The encoded feature vector provides high-dimensional input information with fewer dimensions, making it robust against environmental noise. Furthermore, the above feature vectors can be narrowed down to more effective information for task learning by introducing an attention mechanism (see Section~\ref{sec_attention}).


\item[Time-series Prediction Unit:]\mbox{}\\
%We use a Recurrent Neural Network (RNN) with a recursive connection in the middle layer as the time-series prediction unit. The RNN is a model that reflects temporal dependence, and functions as a learning-based controller that takes the current robot state as input and outputs the next robot state. The RNN takes the feature vector obtained by the sensor feature extractor and current robot state (joint angle torque value, gripper value, etc.) as input and outputs the next robot state, and learns to reduce the error between the output and correct next robot state. At this time, it is possible to construct a robust model by adjusting the ratio of how much of the model output from the previous time is used as input. It is also empirically known that predicting sensor values as a subtask contributes to the predictive performance of the model.
We used a Recurrent Neural Network (RNN) with a recursive connection in the middle layer as the time-series prediction unit. An RNN model reflects temporal dependence and functions as a learning-based controller that inputs the current robot state and outputs the next robot state. The RNN takes the feature vector obtained by the sensor feature extractor and the current robot state (joint angle torque value, gripper value, etc.) as input, outputs the next robot state, and learns to reduce the error between the output and the correct next robot state. A robust model can be constructed by adjusting the ratio of the model output to the previous time as the input. It is also empirically known that predicting sensor values as a subtask contributes to the predictive performance of a model.


%In addition to the above, the RNN can retain various dynamics obtained from the learning process of time-series data as unique internal states~\cite{Tani1998}. When the prediction error increases during robot motion generation, the RNN has the property of switching from the learned dynamics to minimize the prediction error (converge to the appropriate dynamics), and this property can be used to implement the predictive coding strategies described in the previous section. Furthermore, by representing long-horizon tasks as a combination of dynamics embedded in the RNN, it is possible to realize a composite generation of motions. The above details will be described in Section~\ref{sec_case_studies}.
Additionally, RNN can retain various dynamics obtained from the time series data learning process as unique internal states~\cite{Tani1998}. When the prediction error increases during robot motion generation, the RNN can switch from learned dynamics to minimize the prediction error (converging to the appropriate dynamics). This property can be used to implement the predictive coding strategies described in the previous section. Furthermore, the composite generation of motions can be realized by representing long-horizon tasks as a combination of dynamics embedded in the RNN. The details are presented in Section~\ref{sec_case_studies}.
\end{description}


\subsubsection{Robot Motion Generation based on Minimizing Prediction Errors}
\label{sec_generation}

%While the robot performs a task, it sequentially acquires sensor information. By performing forward computation of the model at each step, the RNN predicts the robot's state at the next time step from the input and internal context information. Then, the robot's joints are controlled using the RNN prediction as the target state. By repeating the above process online, the RNN predicts the robot's sensorimotor behavior while sequentially changing the state of each neuron in the context layer. Since our entire system is composed of light models, one of the advantages is that the calculation time and cost required for motion generation are low.
As the robot performs a task, it sequentially acquires sensor information. The RNN predicts the robot's state at the next time step by performing forward computation of the model using the input and internal context information at each step. Subsequently, the robot's joints were controlled using the RNN prediction as the target state. The RNN predicts the robot's sensorimotor behavior while sequentially changing the state of each neuron in the context layer by repeating the above process online. One of its advantages is that the calculation time and cost required for motion generation are low because our entire system is composed of light models.


%During online motion generation, the manipulated object, robot arm, sensor, and the DPL model engage each other. And, the perceptual predictions (posterior beliefs) of the model are corrected in real-time. As a result of the motion modification, the sensory input is also modified, reducing the model's prediction error. At this time, it can be observed that small trajectory corrections and operation switching are occurring through trial and error~\cite{Ichiwara2022contact,Koma2016ral}. This means transitions to other behavior attractors occur due to attraction based on prediction errors. What is important here is that trial and error is not a random process, but a correction of predicted behavior based on experience. We use the above characteristics to accomplish high-dimensional, complex, and long-horizon tasks.
During online motion generation, the manipulated object, robot arm, sensor, and DPL model engage with each other. The model's perceptual predictions (posterior beliefs) were corrected in real-time. As a result of the motion modification, the sensory input was also modified, thereby reducing the model's prediction error. It can be observed that small trajectory corrections and operation switching occur through trial and error~\cite{Ichiwara2022contact,Koma2016ral}. These observations imply that transitions to other behavior attractors occur because of attractions based on prediction errors. What is important here is that trial and error are not random processes but corrections of predicted behavior based on experience. We use these characteristics to accomplish high-dimensional, complex, and long-horizon tasks.


\subsubsection{Hierarchical Structure with Different Time-scales}
\label{sec_timescale}

%%We discuss in more detail the requirements of the RNN in the DPL concept.
%Hierarchy that considers multiple time scales of the system is also important to predictive learning. In other words, it is a system in which a model that makes predictions on a fast timescale similar to the real world and a model that makes predictions on a much slower time scale exist in parallel. We use multiple-time scale RNN (MTRNN~\cite{Yamashita2008emergence}) and LSTM~\cite{Hochreiter1997nc} in the time-series prediction unit and minimize the overall prediction error by an interaction between layers. In the lower layer, which is close to sensory input, the model predicts fast (short time constant) dynamics, and in the upper layer, the model predicts slow (long time constant) dynamics.
A hierarchy that considers multiple time scales of a system is also important for predictive learning. In other words, it is a system in which a model that makes predictions on a fast time scale similar to the real world and a model that makes predictions on a much slower time scale exist in parallel. We used a multiple-time-scale RNN (MTRNN~\cite{Yamashita2008emergence}) and LSTM~\cite{Hochreiter1997nc} in the time-series prediction unit and minimized the overall prediction error by the interaction between layers. In the lower layer, which is close to the sensory input, the model predicts fast (short time constant) dynamics, whereas in the upper layer, the model predicts slow (long time constant) dynamics.


%In the above model, sensory prediction errors are absorbed by changes in the state of the lower layer. At this time, the upper layer generates predictions based on long-term intentions and provides the lower layer with constraints to maintain this state. When a sensory prediction error can no longer be absorbed by the lower layer, the error affects the upper layer, causing a state change (change of intention) in the upper layer. This structure in which the upper and lower layers interact with each other enables consistent yet adaptive behavior generation~\cite{Suzuki2018ral,Saito2021}.
In the above model, the sensory prediction errors were absorbed by changes in the state of the lower layer. The upper layer generates predictions based on long-term intentions, giving the lower layer constraints to maintain this state. When the lower layer can no longer absorb a sensory prediction error, it affects the upper layer, causing a state change (change in intention) in the upper layer. This structure, in which the upper and lower layers interact, enables consistent yet adaptive behavior generation~\cite{Suzuki2018ral,Saito2021}.


\subsubsection{Bottom-up and Top-down Attention}
\label{sec_attention}

%Next, we describe the attention mechanism involved in modifying prediction accuracy. In general, attention mechanisms efficiently allocate a model's resources by concentrating the processing area on a specific target. In robot tasks, it is possible to provide robustness against sudden changes in lighting conditions, placement of multiple objects, and movement of target objects that cause noise. Usually, attention in image processing is generated bottom-up directly from the current input image~\cite{Vaswani2017nips}. For example, methods include emphasizing features such as brightness, color, direction, and orientation, or generating noteworthy keys, queries, and values in self-attention.
Next, we describe the attention mechanism involved in modifying prediction accuracy. Generally, attention mechanisms efficiently allocate a model's resources by concentrating the processing area on a specific target. In robotic tasks, it is possible to provide robustness against sudden changes in lighting conditions, placement of multiple objects, and movement of target objects that cause noise. Typically, attention in image processing is generated bottom-up directly from the current input images~\cite{Vaswani2017nips}. For example, methods emphasize brightness, color, direction, and orientation or generate noteworthy keys, queries, and values in self-attention.


%In the DPL framework, we are conscious of guiding the attention area by prediction. By predicting multiple attention points within an image by using spatial softmax~\cite{finn2016deep}, we limit the areas that can be attended to in the image and have no choice but to ignore other areas. We use the attention point and sensorimotor prediction error as a loss function, and as a result, the attention point is self-organized in that it changes depending on one's actions. In addition to the above bottom-up attention, top-down attention generates attention areas based on the intentions of the upper layers. In \cite{Ito2022integrated}, it is possible to confirm the background other than the object being attended to disappears in a predicted image generated as a posteriori belief from a verbal instruction. Please refer to our related work~\cite{Ito2022integrated,Ichiwara2022contact,Ichiwara2023ral} for details on this attention mechanism.
In the DPL framework, we are conscious of guiding the attention areas through predictions. By predicting multiple attention points within an image using spatial softmax~\cite{finn2016deep}, we limit the areas that can be attended to in the image and have no choice but to ignore others. We use the attention point and sensorimotor prediction error as a loss function. As a result, the attention point is self-organized, changing depending on one's actions. In addition to bottom-up attention, top-down attention generates attention areas based on the intentions of the upper layer. In \cite{Ito2022integrated}, it was possible to confirm that the background, other than the object being attended to, disappeared in a predicted image generated as a posteriori belief from verbal instruction. Please refer to our related work~\cite{Ito2022integrated,Ichiwara2022contact,Ichiwara2023ral} for details regarding this attention mechanism.



\subsection{Discussion of Proposed Concept}
\subsubsection{Perspectives on Embodiment and Dynamic Cognitive Systems}

%In the fields based on physicality, such as behavior-based robotics~\cite{Brooks1991intelligence}, it is believed that intelligent movements emerge from the coupling of the environment, body, and nervous system. In the above fields, the importance of the interaction between the body and the environment is particularly emphasized. For example, in a passive walking mechanism, walking motion is generated through the interaction between a well-designed bipedal link model and a slope~\cite{Garcia1996}. Considering the body's flexibility, mechanisms, and sensor placement, the body is thought to potentially reserve a variety of actions. At this time, it is said that the nervous system only needs to perform the minimum functions, and in the example of walking, it is sufficient for the oscillator system to function.
In fields based on physicality, such as behavior-based robotics~\cite{Brooks1991intelligence}, intelligent movements are believed to emerge from the coupling of the environment, body, and nervous system. The above fields emphasize the importance of the interaction between the body and the environment. For example, in a passive walking mechanism, the walking motion is generated through the interaction between a well-designed bipedal link model and a slope~\cite{Garcia1996}. Considering the body's flexibility, mechanisms, and sensor placement, it is believed to reserve a variety of actions. At this time, the nervous system only needs to perform minimal functions, and in the example of walking, it is sufficient for the oscillator system to function.


%With the above coupling, it is also possible to provide various states of attraction (attractors) to the nervous system. The DPL concept can be seen as a framework in which the above nervous system part is constructed using a deep learning model. For example, there is a discussion of dynamic cognitive systems based on the perspective of dynamical systems proposed by Tani~\cite{Tani1996model}. Here, the brain is not considered simply as an input/output system, but as a predictive system for the world. During interactions through the body, the brain sometimes forms a specific attractor structure with the world, and when the discrepancy becomes large, this connection breaks. After that, it transitions to a new attractor, sometimes through a phase transition, sometimes through a chaotic state. As this process is repeated, cognition changes dynamically without converging to a fixed state. Tani et al. have realized this process using RNN and a real robot system, and have been conducting interesting research for many years in an attempt to understand the brain cognitive process~\cite{tani2016exploring}. Although the above study was proposed before the discussion of predictive coding and the free energy principle, it has many commonalities with our DPL concept.
This coupling can provide various states of attraction (attractors) to the nervous system. The DPL concept can be considered a framework in which the above nervous system part is constructed using a deep learning model. For example, there is a discussion of dynamic cognitive systems based on the perspective of dynamical systems proposed by Tani~\cite{Tani1996model}. Here, the brain is not simply considered an input/output system but a predictive system for the world. During interactions with the body, the brain sometimes forms a specific attractor structure with the world; this connection breaks when this discrepancy becomes large. Subsequently, it transitions to a new attractor, sometimes through a phase transition or a chaotic state. As this process is repeated, cognition changes dynamically without converging to a fixed state. Tani et al. realized this process using RNN and a real robot system and have been conducting interesting research for many years, attempting to understand the brain cognitive process~\cite{tani2016exploring}. Although the above study was proposed before the discussion on predictive coding and the free-energy principle, it has many similarities with our DPL concept.


\subsubsection{Contrasting with Reinforcement Learning}

%Finally, we discuss the main differences between reinforcement learning and predictive coding. The discussion here is just a comparison with simple reinforcement learning. As mentioned in the previous chapter, one of the distinctive charms of reinforcement learning is that learning goals can be set freely (for example, scalar values, etc.). The main objective of predictive coding is to minimize the prediction error when the model connects the body and the world. Therefore, it is difficult in principle to set other learning goals in predictive coding. Therefore, unless special measures are taken, it is difficult in principle to set other learning goals in predictive coding.
Finally, we discuss the main differences between reinforcement learning and predictive coding. The discussion here is a simple comparison with simple reinforcement learning. As mentioned in the previous section, one of the distinctive characteristics of reinforcement learning is that learning goals can be freely set (for example, scalar values, etc.). The main objective of predictive coding is to minimize prediction errors when the model connects the body and the world. Therefore, setting other learning goals for predictive coding is difficult. Therefore, setting other learning goals for predictive coding will be difficult unless special measures are taken.


%In model-based reinforcement learning, a method has been proposed to learn a "world model" by repeating interactions between actions and states~\cite{Ha2018world}. The method learns actions and visual field images as a sequence of abstracted latent spaces. This is similar in concept to predictive coding in that it targets the world as a prediction target. After learning has progressed sufficiently, the optimal behavior based on the world model can be obtained by searching through trial and error for the behavioral policy that yields the maximum reward from the model.
In model-based reinforcement learning, a method was proposed for learning a real-world model" by repeating the interactions between actions and states~\cite{Ha2018world}. This method learns actions and visual field images as a sequence of abstracted latent spaces. The learning process is similar to predictive coding in that it targets the world as a prediction target. After learning has progressed sufficiently, the optimal behavior based on the world model can be obtained by searching through trial and error for a behavioral policy that yields the maximum reward from the model.


%On the other hand, predictive coding assumes that the world is open and that the environment model learning is constantly progressing. While learning is insufficient, actions are generated that are expected to reduce future prediction errors (expected free energy). By repeating the experience of reaching the goal state, the sequence of actions becomes a habit as a prior belief. Thus, actions with high prediction accuracy are selected. In this process of habituation, the generated action sequence (virtual expectations) can be thought of as trial and error. In other words, the behavior learned in this framework is not necessarily optimal for a given objective index. The search that accompanies learning has the aspect of not only optimization but also diversification of possible adaptations. The paradox of trying to adapt to the environment but not being able to do so completely leads to the exploration of various actions.
In contrast, predictive coding assumes that the world is open and that environmental model learning is constantly progressing. When the learning is insufficient, actions expected to reduce future prediction errors (expected free energy) are generated. By repeating the experience of reaching the goal state, the sequence of actions becomes a habit, as a prior belief. Thus, actions with high prediction accuracy were selected. The generated action sequence (virtual expectations) can be considered trial and error in this habituation process. In other words, the behavior learned in this framework is not necessarily optimal for a given objective index. The search that accompanies learning involves optimization and diversification of possible adaptations. The paradox of trying to adapt to the environment but not being able to do so leads to the exploration of various actions.


